\scanpage{30}
\uline{Reformuler le th.}

\uline{Th} $X$ propre et plat sur $S$, $H^0(X_s,\underline{\mathrm O}_{X_s})\xleftarrow{\sim}k(s)$ pour tt $s\in S$, ($S$ loc. \unsure{nœth|noeth.}), $\varepsilon_X$ une section de $X/S$, \deleted{\unsure{$\mathrm{Hom}(X$|$\mathrm{Hom}(X,$}} $G$ un foncteur sur $(\unsure{\mathrm{Sch}|\mathrm M})_{/S}^{\circ}$ à valeurs dans la catégorie des groupes. \deleted{On suppose donné un} \inserted{$f:X\longrightarrow G$ un $S$-hom pointé.} \deleted{On suppose donné un hom}

(i) Si $S_0$ est un sous-préschéma \deleted{\uline{\unsure{réduit}}} de $S$ ayant $=$ ensemble sous-jacent, et si \deleted{$G$ est} l'hom $f_0:X_0\longrightarrow G_0$ est trivial, alors $f$ l'est [ pourvu que \unsure{--- petite condition|--- petite cartésienne|--- petite catégorie} )

(ii) si de plus $G$ commute avec $\varprojlim$ d'anneaux artiniens, aux $\varinjlim$ d'anneaux noeth. \inserted{[cas de la localisation]}, et \deleted{\unsure{les}} les morphismes $\operatorname{Spec}(\widehat A)\longrightarrow\operatorname{Spec}(A)$ sont des morphismes de $G$-\unsure{0|O}-descente, \uline{alors} si $f_{\unsure{\times s|s}}=f\times_S k(s)$ est trivial, $f$ est trivial dans un voisinage

(iii) Si $G$ est séparé sur $S$, alors sous les conditions de (ii), \deleted{\unsure{on}} $f$ est trivial sur la composante connexe de $\unsure{e|s}$

En d'autres termes, $\underline{\mathrm{Hom}}_S^{\bullet}(X,G)$ est \uline{non ramifié} sur $\unsure{S|?}$
% archive sheet number: 29
